\section{Methodology}
\label{sec:method}

The divergence mask of Eq.~(\ref{eq:dppo}) upgrades the trust-region
\emph{proximity} criterion from a sampled-ratio proxy to a distributional
divergence, but leaves the
\emph{direction} criterion $\sign(\Adv_t(r_t-1))$ untouched. This section makes the
direction criterion consistent with the divergence-based proximity criterion. The
principle is to ask whether the next gradient step would actually increase the
divergence (\S\ref{sec:method-principle}).
We answer it with the divergence's own directional derivative, which we derive in
closed form (\S\ref{sec:method-derivation})
and then recover from the truncated top-$K$ statistics a rollout reports
(\S\ref{sec:method-estimators}). The result is a drop-in replacement for the direction
criterion, together with an analysis of when it agrees with the ratio-based one and when
it does not (\S\ref{sec:method-mask}).

\subsection{From ``will the ratio move away from one?'' to ``will the divergence grow?''}
\label{sec:method-principle}

The direction criterion determines whether keeping a token update would drive
the policy \emph{further} outside the trust region.
Under a ratio-based trust region, this question is about $|r_t-1|$.
The criterion $\sign(\Adv_t(r_t-1))$ answers it exactly because it is computed from the same ratio that defines the trust region.
Its sign therefore reports whether the update increases $|r_t-1|$.
Under a divergence-based trust region, the proximity criterion is instead defined by $\Dtok_t$, which aggregates the gap between the two policies over the entire vocabulary, while $\sign(\Adv_t(r_t-1))$ still looks only at the single sampled token. 
Its sign can therefore disagree with the sign of the actual change in $\Dtok_t$, 
so the ratio-based direction criterion no longer reliably indicates whether the update drives the policy further outside the trust region.

We therefore define the direction criterion directly in terms of $\Dtok_t$.
It asks whether the next gradient step increases $\Dtok_t$.
Let $\tpi_\eta$ denote the training policy after moving its logits along the
surrogate gradient by an effective step $\eta\ge0$.
Expanding the divergence along this step,
\begin{equation}
\label{eq:taylor-def}
\Dtok_t(\tpi_\eta) \;=\; \Dtok_t \;+\; \eta\cdot \frac{d}{d\eta}\Dtok_t(\tpi_\eta)\Big|_{\eta=0} \;+\; O(\eta^2),
\end{equation}
the sign of the first-order term determines the local trend of $\Dtok_t$.
A positive sign means the update increases $\Dtok_t$, and a negative sign means
it decreases $\Dtok_t$.
A token already outside the trust region should therefore be blocked when this
term is positive, since the update would push $\Dtok_t$ still higher, and kept
when it is negative, since the update would pull $\Dtok_t$ back down.
We replace the ratio-based direction criterion with the sign of this first-order
change, the divergence's own response to the update rather than a single-sample
proxy for it.

\subsection{The directional derivative of the divergence}
\label{sec:method-derivation}

We work at a single token and drop the index $t$ for brevity. Let
$\tpi_i = e^{z_i}/\sum_j e^{z_j}$ be the training policy over the vocabulary,
with logits $z_i$, and let $k$ be the sampled action (the token $y_t$ of the
underlying step). The per-token surrogate
objective $r\,\Adv$ has gradient $\Adv\,r\,\nabla_z\log\tpi_k$ with respect to the
logits, where $\nabla_{z_i}\log\tpi_k=\1[i=k]-\tpi_i$. A single gradient step therefore moves the
logits along the direction $\1[i=k]-\tpi_i$, scaled by $\Adv\,r$.
Since $r>0$, the \emph{sign} of the step is $\sign(\Adv)$.
We compute the directional derivative for a
unit step in $\bm{v}_i=\1[i=k]-\tpi_i$ and recombine the advantage sign afterwards.

\begin{proposition}[Directional derivative of $\KL(\bmu\,\|\,\tpi)$ along a policy-gradient step]
\label{prop:dir-deriv}
Let $\tpi_\eta$ be the softmax policy whose logits are $z_i+\eta\,(\1[i=k]-\tpi_i)$.
Then
\begin{equation}
\label{eq:taylor-coeff}
\Taylor \;\equiv\; \frac{d}{d\eta}\,\KL(\bmu\,\|\,\tpi_\eta)\Big|_{\eta=0}
\;=\; (\tpi_k-\bmu_k) \;+\; \sum_i \tpi_i\,(\bmu_i-\tpi_i).
\end{equation}
\end{proposition}

We denote this directional derivative $\Taylor$ throughout and defer its derivation to
Appendix~\ref{app:derivations}.

\noindent\textbf{What the ratio-based direction criterion misses.}
The coefficient $\Taylor$ decomposes into two terms with distinct scope,
\begin{equation}
\label{eq:decomp}
\Taylor \;=\; \underbrace{(\tpi_k-\bmu_k)}_{\text{local: sampled token } k}
\;+\; \underbrace{\sum_i \tpi_i(\bmu_i-\tpi_i)}_{\text{global: normalization coupling}},
\end{equation}
and the split is exactly the boundary between what a single sample observes and what
it cannot. The local term is the probability gap at the sampled token. Since
$r-1=(\tpi_k-\bmu_k)/\bmu_k$ with $\bmu_k>0$,
\begin{equation}
\label{eq:local-is-ratio}
\sign(\tpi_k-\bmu_k)=\sign(r-1),
\end{equation}
so the ratio-based direction criterion is precisely the sign of the local term.
The decomposition shows exactly what this criterion misses.
Changing the sampled probability necessarily changes the rest of the distribution
through the softmax normalizer, creating the global term $\E_{i\sim\tpi}[\bmu_i-\tpi_i]$ that no
single sampled ratio can reveal.
The ratio-based direction criterion is therefore equivalent to using only the
local term and ignoring this coupling.
The divergence-based criterion keeps both terms, which is why its sign can track
the true change in $\Dtok_t$ where the ratio sign cannot.

\begin{remark}[Step-size independence]
\label{rem:step-size}
The actual logit step scales the unit direction $\bm v$ by $\eta\,\Adv\,r$,
where $\eta>0$ is the effective step size and $r>0$ is the importance ratio.
Thus the first-order change of the divergence has sign
$\sign(\Adv\cdot\Taylor)$.
The mask only needs this sign.
It does not need to estimate $\eta$ or the magnitude of $r$.
By contrast, predicting the post-update value $\Dtok_t(\tpi_\eta)$ would require
those quantities and the higher-order $O(\eta^2)$ terms.
The sign prediction is reliable when the first-order term dominates the
higher-order remainder.
\end{remark}

\subsection{Estimating the directional derivative under a truncated vocabulary}
\label{sec:method-estimators}

Computing $\Taylor$ exactly would require the full vocabulary distribution, whereas a
production rollout stores only the top-$K$ probabilities and the leftover tail mass.
The shortfall is minor, because the sampled token and the retained top-$K$ entries are
observed directly and the only missing ingredient is the contribution of the unseen
tail. Estimating $\Taylor$ therefore reduces to modeling that tail, and we
consider two estimators that differ only in how they do so.

\noindent\textbf{Aggregated-tail estimator.}
The simplest choice mirrors the bucketing already used by the top-$K$ divergence in
Eq.~(\ref{eq:topk-kl}) and collapses the entire tail into one aggregated token of mass
$\tpi_{\mathrm{tail}}$ (resp.\ $\bmu_{\mathrm{tail}}$). The coefficient
becomes
\begin{equation}
\label{eq:agg-tail}
\Taylor^{\mathrm{agg}}=(\tpi_k-\bmu_k)
+\sum_{i\in\Keep}\tpi_i(\bmu_i-\tpi_i)
+\tpi_{\mathrm{tail}}\big(\bmu_{\mathrm{tail}}-\tpi_{\mathrm{tail}}\big).
\end{equation}
This single-bucket approximation can overestimate the tail contribution by
concentrating all residual mass on one aggregate token.

\noindent\textbf{Uniform-tail estimator.}
A more faithful model spreads the residual mass evenly over the $n-m$ unseen
tokens, taking $\tpi_j\approx\tpi_{\mathrm{tail}}/(n-m)$ and
$\bmu_j\approx\bmu_{\mathrm{tail}}/(n-m)$
for $j\notin\Keep$ (with $m=|\Keep|$ and $n$ the vocabulary size). Substituting into
the tail sums gives
\begin{equation}
\label{eq:uni-tail}
\Taylor^{\mathrm{uni}}=(\tpi_k-\bmu_k)
+\sum_{i\in\Keep}\tpi_i(\bmu_i-\tpi_i)
+\frac{\tpi_{\mathrm{tail}}\big(\bmu_{\mathrm{tail}}-\tpi_{\mathrm{tail}}\big)}{n-m},
\end{equation}
the same tail term as Eq.~(\ref{eq:agg-tail}) but suppressed by $1/(n-m)$. Since the
vocabulary is large this term is negligible, in contrast to the exaggerated estimate of
the single-bucket model.

\noindent\textbf{The gap between the two estimators.}
Since the two estimators share all retained-token terms and differ only in their
tail contribution, their difference is
\begin{equation}
\label{eq:tail-gap}
\Taylor^{\mathrm{agg}}-\Taylor^{\mathrm{uni}}
=\Big(1-\tfrac{1}{n-m}\Big)\,\tpi_{\mathrm{tail}}\,(\bmu_{\mathrm{tail}}-\tpi_{\mathrm{tail}})
\;\approx\; \tpi_{\mathrm{tail}}\,(\bmu_{\mathrm{tail}}-\tpi_{\mathrm{tail}}),
\end{equation}
where the approximation uses $n-m\gg1$.
In LLMs, the vocabulary size $n$ is typically on the order of $10^5$, while
$m=|\Keep|$ is typically at most $20$.
Thus the gap is small whenever the retained top-$K$ tokens capture most of the
probability mass.
In this regime, the aggregated-tail and uniform-tail masks should behave nearly
identically.

Both estimators are computationally lightweight. Each reads only the top-$K$
probabilities and tail masses already returned by the rollout, requiring no
additional forward or backward pass and no finite-difference evaluation.
Appendix~\ref{app:binary-kl-equivalence} shows that, under DPPO's binary KL
approximation, the aggregated-tail divergence-based direction criterion reduces exactly to the
ratio-based direction criterion.
Appendix~\ref{app:topk-tv-predictive} gives the analogous divergence-based direction criterion
when the proximity criterion is top-$K$ TV.

\subsection{Policy optimization with the predictive divergence mask}
\label{sec:method-mask}

Replacing the ratio-based direction criterion in Eq.~(\ref{eq:dppo}) with the
divergence-based direction criterion built from the directional derivative $\Taylor_t$ gives
the \emph{predictive divergence mask}:
\begin{equation}
\label{eq:predictive-mask}
M_t^{\mathrm{pred}} =
\begin{cases}
0, & \sign\big(\Adv_t\cdot\Taylor_t\big)>0 \;\text{ and }\; \Dtok_t>\delta,\\[2pt]
1, & \text{otherwise,}
\end{cases}
\end{equation}
where $\Dtok_t=\KLtopk(\bmu\,\|\,\tpi)$ and $\Taylor_t$ is the directional derivative of
Eq.~(\ref{eq:taylor-coeff}) (estimated by either the aggregated-tail or uniform-tail form). The mask reuses only
quantities already on hand during training, namely $\Adv_t$, the top-$K$ probabilities,
and the divergence, and introduces no new hyperparameter beyond the existing
threshold $\delta$. The resulting objective is the masked policy gradient
\begin{equation}
\label{eq:predictive-obj}
\nabla_\theta L^{\mathrm{pred}}(\tpi)
= \E_{y\sim\bmu}\!\left[\sum_{t=1}^{|y|}
M_t^{\mathrm{pred}}\; r_t\,\Adv_t\;\nabla_\theta\log\tpi(y_t\mid s_t)\right],
\end{equation}
which is identical to the DPPO-TopK-KL update of Eq.~(\ref{eq:dppo}) except that its direction criterion follows the divergence's first-order change instead of the sampled ratio.

\begin{remark}[Local nature of the prediction]
The predictive mask is a token-level first-order approximation.
It estimates whether one token's gradient contribution would locally increase or
decrease that token's divergence.
The actual parameter update aggregates gradients from all tokens in the batch, so
the realized divergence change at one token can also be affected by other tokens'
updates.
Thus the predicted sign should be viewed as a local estimate of the divergence
trend, not an exact prediction of the post-update divergence.
\end{remark}

\noindent\textbf{Agreement and disagreement of the two criteria.}
The decomposition in Eq.~(\ref{eq:decomp}) can also be read as centering the
sampled-token gap by the policy-weighted average gap:
\[
\Taylor
=
(\tpi_k-\bmu_k)-\E_{i\sim\tpi}[\tpi_i-\bmu_i].
\]
Thus the global term changes the mask decision only if this centering crosses
zero.
Equivalently, disagreement occurs only when the current-policy average gap lies
farther from zero in the same direction as the sampled-token gap:
\begin{equation}
\label{eq:recovery}
0<\tpi_k-\bmu_k<\E_{i\sim\tpi}[\tpi_i-\bmu_i]
\quad\text{or}\quad
0>\tpi_k-\bmu_k>\E_{i\sim\tpi}[\tpi_i-\bmu_i].
\end{equation}
Together with the proximity condition $\Dtok_t>\delta$, these are precisely the
tokens on which the predictive mask differs from the ratio-based mask.
They are also the cases where a single-sample ratio is least informative: the
sampled-token gap points in one direction, but the distribution-wide correction
reverses the predicted change in divergence.
Otherwise, the centered gap keeps the local sign, and the two masks agree.
